
\subsection{Twist-ladder closure convention}
    \label{subsec:twist_ladder_closure_convention}

    \textbf{[ORTHODOX TOPOLOGY].}
    A clasp closure with a twist box containing \(m\) half-twists generates the
    standard twist-knot ladder. In the convention used here,
    \begin{align}
        K_1 &= 3_1, &
        K_2 &= 4_1, &
        K_3 &= 5_2, &
        K_4 &= 6_1,
    \end{align}
    and, for \(m\geq 1\),
    \begin{align}
        c_{\min}(K_m)=m+2,
        \qquad
        \det(K_m)=2m+1 .
        \label{eq:twist_knot_ladder_invariants}
    \end{align}
    Reversing the handedness of the twist box gives the mirror class
    \begin{align}
        K_{-m}=\overline{K_m} .
    \end{align}

    \textbf{[SPECULATIVE SST INTERPRETATION].}
    The twist-ladder construction is retained as a mechanical template for
    R-to-T closure: the integer \(m\) counts localized half-twist content before
    clasp closure, while the sign of \(m\) records chirality. The electron
    candidate corresponds to the minimal chiral closure
    \begin{align}
        m=1 \quad\Longrightarrow\quad K_m=3_1 .
    \end{align}
    This is a generative convention, not yet a first-principles filament theorem.

\subsection{Euler topological preservation and closure discipline}
    \label{subsec:euler_topological_preservation_closure}

    \textbf{[ORTHODOX].}
    In a smooth incompressible inviscid Euler flow without boundaries, forcing,
    viscosity, singular core collapse, or phase-slip defects, vortex lines are
    materially transported by the flow. For
    \begin{align}
        \partial_t\mathbf{v}
        +
        (\mathbf{v}\cdot\nabla)\mathbf{v}
        =
        -\frac{1}{\rhoF}\nabla p,
        \qquad
        \nabla\cdot\mathbf{v}=0,
    \end{align}
    the vorticity \(\boldsymbol{\omega}=\nabla\times\mathbf{v}\) satisfies
    \begin{align}
        \partial_t\boldsymbol{\omega}
        =
        \nabla\times(\mathbf{v}\times\boldsymbol{\omega}),
        \label{eq:euler_vorticity_frozen_in}
    \end{align}
    and Kelvin circulation is conserved on a material loop,
    \begin{align}
        \frac{D}{Dt}\oint_{C(t)}\mathbf{v}\cdot d\boldsymbol{\ell}=0 .
        \label{eq:kelvin_preservation_closure_section}
    \end{align}
    Consequently, linking and knot type of already closed vortex tubes are
    preserved under smooth ideal evolution:
    \begin{align}
        0_1 \not\to 3_1
        \qquad
        \text{under smooth ideal Euler flow.}
        \label{eq:unknot_not_to_trefoil_euler}
    \end{align}

    \textbf{[DERIVED SST CONSTRAINT].}
    Trefoil formation must not be described as ordinary reconnection of an
    already closed vortex tube in the ideal bulk. Trefoil creation is admissible
    only as a non-Euler closure event, boundary event, phase-slip event, or
    nucleation event occurring before the final closed vortex carrier has
    acquired a protected knot type.

    \textbf{[SPECULATIVE SST BRIDGE].}
    The Kairos closure event is therefore modeled as
    \begin{align}
        \gamma_{\rm R}
        +
        \mathcal{B}_{\rm atom}
        \overset{\mathcal{K}_{\rm Kairos}}{\longrightarrow}
        3_1
        +
        \mathcal{R},
        \label{eq:kairos_closure_event}
    \end{align}
    where \(\gamma_{\rm R}\) is an extended R-phase torsion pulse,
    \(\mathcal{B}_{\rm atom}\) is a local boundary or trapping condition, and
    \(\mathcal{R}\) carries recoil and topological compensation. The local
    transition scale is \(a_{\rm core}\), while the post-closure neutral
    circulation envelope is normalized by \(\rc\).
